\begin{proposition}[Proved: modular regular coherence is not necessary in the all-ring problem]
\label{C:noncoherent-positive-boundary}
Let $A$ be any mixed-characteristic discrete valuation ring,
with uniformizer $\pi$, residue field $\kappa$ of characteristic
$p>0$, and fraction field $K$.
Let $x,y$ be algebraically independent over $K$.  In $K(x,y)$ put
\[
 P_n=\pi/x^n,\qquad Y_n=y/x^n,\qquad
 R_n=A[x,P_n,Y_n]_{(x,P_n,Y_n)}\quad(n\geq0),
 \qquad S=\bigcup_{n\geq0}R_n.                     \tag{C.1.NC1}
\]
Here $P_0=\pi$, so
$R_0=A[x,y]_{(\pi,x,y)}$.
Then $S$ is a characteristic-zero normal local domain with
maximal ideal $xS$, residue field $\kappa$, and
$K_0(S)=\mathbb Z$.  The map $\vartheta_S$ is an isomorphism
in every degree, although $S$ is not coherent and is not
a valuation domain.  In particular its localization at its
positive-residue-characteristic closed point is not
regular coherent.

Thus regular coherence at all positive-residue-characteristic
points cannot be a necessary condition for the intrinsic
all-ring C.1 criterion.  This statement does not decide
necessity of modular regularity within the class of coherent
rings.  Its positive comparison assertion is an application
of the already proved filtered-colimit theorem, not a new
family beyond that theorem.
\end{proposition}
\begin{proof}
\emph{The regular local stages.}
For $n\geq1$, the indicated subring of $K(x,y)$ has presentation
\[
 A[x,P_n,Y_n]\simeq
 A[X,P,Y]/(X^nP-\pi),                           \tag{C.1.NC2}
\]
where $X$ maps to $x$ and $Y$ to $y/x^n$.
Indeed $X^nP-\pi$ is primitive and irreducible as a
polynomial in $P$ over the UFD $A[X,Y]$: its two
coefficients are relatively prime, and it has degree one
over the fraction field.
It is therefore prime in $A[X,P,Y]$.
The quotient embeds after inverting $X$ in
$A[X,X^{-1},Y]$, where $P=\pi/X^n$; replacing $Y$ by
$y/x^n$ is an embedding into $K(x,y)$.
This proves the presentation and that its displayed ideal
$(X,P,Y)$ is maximal with residue field $\kappa$.

The ambient local ring
$A[X,P,Y]_{(\pi,X,P,Y)}$ is regular of dimension four.
Its relation $X^nP-\pi$ has nonzero image $-\pi$ in the
cotangent space, since $n\geq1$.
Quotienting by it therefore gives a regular local ring
$R_n$ of dimension three and regular parameters
$x,P_n,Y_n$.
The same conclusion for $R_0$ is immediate from its
polynomial presentation.  In particular every $R_n$ is
a nonzero characteristic-zero regular noetherian local ring.

\emph{The local union and its comparison map.}
The substitutions
\[
 x\longmapsto x,\qquad
 P_n\longmapsto xP_{n+1},\qquad
 Y_n\longmapsto xY_{n+1}                         \tag{C.1.NC3}
\]
give injective local homomorphisms $R_n\to R_{n+1}$:
they are the inclusions into $K(x,y)$, and on residue
fields they are the identity of $\kappa$.
The union is consequently a local domain whose maximal
ideal is the union of the stage maximal ideals and whose
residue field is $\kappa$.
Every generator $x,P_n,Y_n$ of a stage maximal ideal belongs
to $xR_{n+1}$.  Hence the maximal ideal of $S$ is exactly
$xS$; in particular $x$ is not a unit.
Finite projective modules over a local ring are free, so
the rank identifies $K_0(S)$ with $\mathbb Z$.

Normality can also be checked directly, without a
noetherian hypothesis on the union.  If a fraction of $S$
is integral over $S$, choose one stage containing its
numerator, its nonzero denominator and all coefficients
of a monic integral equation.
This fraction belongs to the fraction field of that
normal stage and is integral over it, so belongs to the
stage.  Thus $S$ is integrally closed.

Theorem~\ref{C:regular_noetherian} gives that
$\vartheta_{R_n}$ is an isomorphism for every $n$.
Proposition~\ref{C:filtered_colimits} now gives the same
assertion for $S$.  Its finite-matrix descent proof applies
to these nonflat transition maps: admissible sequences
are split on their underlying projective modules, and
no flatness assumption is present in that proposition.

\emph{A non-finitely-generated intersection of principal ideals.}
We claim
\[
 I:=\pi S\cap yS
       =\bigcup_{n\geq0}(\pi y/x^n)R_n.         \tag{C.1.NC4}
\]
Every term on the right belongs to the intersection,
since $y/x^n$ and $\pi/x^n$ belong to $R_n$.
Conversely let $a\in I$, and write $a=\pi b=yc$
with $b,c\in S$.
Choose $n$ such that $a,b,c\in R_n$.  Cancelling the
nonzero $x^n$ gives
$P_n b=Y_n c$ in $R_n$.
The ideal $P_nR_n$ is prime, and $Y_n$ does not belong
to it, since
\[
 R_n/P_nR_n\simeq\kappa[x,Y_n]_{(x,Y_n)}.
\]
This also holds for $n=0$, where $P_0=\pi$.
Thus $c=P_nd$ for $d\in R_n$; cancelling $P_n$
then gives $b=Y_nd$ and
$a=(\pi y/x^n)d$.
This proves (C.1.NC4).
The terms on the right are increasing because
$\pi y/x^n=x(\pi y/x^{n+1})$.

Dividing such an expression for $a$ by $x$ gives
$(\pi y/x^{n+1})d\in I$.
Consequently $I=xI$.  Since $I$ contains the nonzero
element $\pi y$ and $x$ belongs to the maximal ideal of
$S$, Nakayama's lemma implies that $I$ is not finitely
generated.
The intersection of two finitely generated ideals in a
coherent ring is finitely generated.  Applied to the
principal ideals $\pi S$ and $yS$, this proves that $S$
is not coherent.  In particular it is not noetherian.

Finally $\pi/y=P_n/Y_n$ cannot belong to any $R_n$,
because $Y_nR_n$ is prime and does not contain $P_n$.
Likewise $y/\pi$ cannot belong to any $R_n$, by the
primeness calculation just used.
Thus neither ratio belongs to $S$, proving that $S$ is
not a valuation domain.
The closed-point localization is $S$ itself and has
residue characteristic $p$, so it is not regular coherent.
This gives exactly the stated boundary on necessity.
\end{proof}

\noindent\emph{Source and scope.}
The rings in (C.1.NC1) form a sequence of local quadratic
transforms, so their union is a quadratic Shannon
extension.  Noncoherence of a nonvaluation quadratic
Shannon extension is a known general result:
\cite[Theorem~6.2, especially its nonarchimedean
Case~1]{GuerrieriHeinzerOlberdingToeniskoetter2017}.
The concrete principal-ideal calculation above supplies
the entire needed proof in this mixed-characteristic
family.  No novelty claim is made for the commutative
algebra example, and no finite computation is substituted
for its non-finite-generation argument.
